\section{A sheet with finitely many parameters}
\label{sec:lvsheet}

\subsection{Positivity as a box}

The unknown is now a positive function of two variables --- an
infinite-dimensional object with a one-sided constraint, which sounds
harder than anything fitted so far in this book.  The finite version
is built from the simplest ingredient available: straight-line
interpolation.  Picture the smallest example that shows everything.
Take two maturities and three strikes --- a grid of six points --- and
write a positive number on each.  Split each of the two rectangular
cells along a diagonal, giving four triangles.  Inside each triangle,
define the field as the unique plane through its three corner values.
The result is a continuous surface made of flat triangular facets,
passing through the six numbers --- and those six numbers are the
model's only parameters.  The general object just replaces six by a
few hundred.

Two names before the formal definition.  A \emph{tensor grid} is a
rectangular grid: pick a list of maturities and a list of strikes, and
take every combination as a vertex.  The \emph{hat function} of a
vertex is the tent that rises to $1$ there, slopes to $0$ at every
neighboring vertex, and stays $0$ beyond; on each triangle it is the
affine function with those corner values.  Its values at a point are
also called the point's \emph{barycentric weights} --- the three
nonnegative numbers, summing to one, that locate the point inside its
triangle.

\begin{definition}[$P_1$ local-variance sheet]
\label{def:lvsheet}
Fix vertices $(\tau_i,y_j)$, $i=1..n_\tau$, $j=1..n_y$, on a tensor grid,
triangulated by splitting each rectangular cell along one diagonal
(\cref{rem:lvdiag}).  The local variance is the continuous
piecewise-affine interpolant
\begin{equation}
v_\theta(\tau,y)=\sum_{\ell=1}^{n_\theta}\theta_\ell\,
\mathcal H_\ell(\tau,y),\qquad n_\theta=n_\tau n_y,
\label{eq:lvp1}
\end{equation}
where $\mathcal H_\ell$ is the hat function of vertex $\ell$ ---
barycentric on each triangle, equal to $1$ at its own vertex and $0$ at
every other --- and the parameters $\theta_\ell=v(\tau_i,y_j)>0$ are the
\emph{nodal local variances}, indexed flat as $\theta_\ell$ or by row and
column as $\theta_{i,j}$, whichever is convenient.
\end{definition}

With the names in place, read \eqref{eq:lvp1} as plainly as it
deserves: the sum simply interpolates the nodal values.  At any point
of the grid, only the three hats of the surrounding triangle are
nonzero; their values there are the point's barycentric weights; and
$v_\theta$ at that point is the corresponding weighted average of the
three corner parameters.

Why affine pieces, and why triangles?  Because a weighted average with
nonnegative weights cannot leave the range of the numbers it
averages.  That one-line fact is the model's entire validity theory:

\begin{proposition}[Nodal bounds are surface bounds]
\label{prop:lvnodal}
If $v_{\rm lo}\le\theta_\ell\le v_{\rm hi}$ for all $\ell$, then
$v_{\rm lo}\le v_\theta(\tau,y)\le v_{\rm hi}$ everywhere on the
triangulation's hull.
\end{proposition}

\begin{proof}
At any point of the hull, the only nonzero hat functions are the three
belonging to the surrounding triangle, and their values there are
nonnegative and sum to one --- they are the point's barycentric
weights.  Hence
\[
v_\theta(\tau,y)
=\sum_{\ell}\theta_\ell\,\mathcal H_\ell(\tau,y)
\;\le\;v_{\rm hi}\sum_{\ell}\mathcal H_\ell(\tau,y)
=v_{\rm hi},
\]
and the lower bound is the same computation with $v_{\rm lo}$.
\end{proof}

Positivity of a surface has been reduced to positivity of finitely
many numbers: keep every $\theta_\ell$ inside the box
$[v_{\rm lo},v_{\rm hi}]$, and the whole field stays inside it too,
everywhere on the hull --- the region the triangles cover.  A
box is the one constraint that least-squares solvers handle natively
--- clip each coordinate --- with no penalty terms and no feasibility
questions.  This is the same chart construction, one dimension up: as with the
LQD chart of Chapter~2 and the structural chart of
\cref{sec:svistructural}, the coordinates are chosen so that every
point of the parameter set is a valid object, and the optimizer never
learns that a constraint existed.  Combined with
\cref{prop:lvpropagate}, the chain of guarantees is complete before
any data arrives:
\[
\text{box}\;\Longrightarrow\;\text{positive field}
\;\Longrightarrow\;\text{butterfly-free, calendar-ordered surface.}
\]
Every iterate of every fit in this chapter is an admissible diffusion.

\begin{figure}[htbp]
\centering
\paperfig{0.9\textwidth}{fig_lv_sheet.pdf}
\caption{The model, drawn: the calibrated SPY local-volatility sheet of
\cref{sec:lvexamples} (\MacLvSheetRows\ maturity rows $\times$
\MacLvSheetCols\ strike columns, \MacLvSheetVtx\ vertices), with the
triangulation the pricing basis actually uses.  Every dot is one
calibration parameter; every triangle is a region where local variance
is affine in $(\tau,y)$.}
\label{fig:lvsheet}
\end{figure}

\Cref{fig:lvsheet} shows the object this chapter will calibrate: the
SPY sheet of \cref{sec:lvexamples}, with \MacLvSheetRows\ maturity
rows, \MacLvSheetCols\ strike columns, and therefore \MacLvSheetVtx\
vertices, drawn with the very triangulation the pricing uses.  Read it
as a list of numbers, because that is all it is.  Every dot is one
calibration parameter; every triangle is a region where local variance
is an affine function of $(\tau,y)$.  The tall put-side wall, the
sharp structure at the short end, and the flat call-side plain are
nothing but nodal values.  Notice also that the strike columns crowd
toward the money at short maturities; that is deliberate grid design,
and \cref{sec:lvinfo} explains it.

\begin{figure}[htbp]
\centering
\paperfig{\textwidth}{fig_lv_basis.pdf}
\caption{The anatomy of the sheet, on a small demonstration grid.
(a)~The triangulated vertex grid; the shaded star of triangles is the
support of one vertex's hat function.  The cell diagonals are not
uniformly oriented; that is real (\cref{rem:lvdiag}).  (b)~A strike
cross-section along a vertex row: each hat is $1$ at its own vertex,
$0$ at every other, affine between, and they sum to $1$ identically.}
\label{fig:lvbasis}
\end{figure}

\Cref{fig:lvbasis} takes the construction apart on a small
demonstration grid.  In panel~(a), the triangulated vertex grid; the
shaded star of triangles around one vertex is the support of that
vertex's hat function --- the entire region of the surface its
parameter can influence.  A nodal variance is a local dial: turn it,
and only the prices of options whose diffusion passes through that
star can move.  (The cell diagonals in the panel do not all point the
same way; that is real, and \cref{rem:lvdiag} explains it.)  In
panel~(b), a cross-section along one vertex row: each hat rises to $1$
at its own vertex, is $0$ at every other, and the hats sum to $1$
identically --- the partition of unity behind the weighted-average
argument of \cref{prop:lvnodal}.

\begin{remark}[Which diagonal, and does it matter]
\label{rem:lvdiag}
On a tensor grid every cell is a rectangle, and a rectangle's four
corners lie on one circle --- exactly the tie case of the standard
(Delaunay) criterion for well-shaped triangles, which prefers
triangulations in which no vertex falls strictly inside another
triangle's circumscribed circle.  The triangulation routine therefore breaks the tie cell by
cell, deterministically but by no designed rule, and the figures draw
the resulting patchwork faithfully.  The choice is second-order small:
the two candidate interpolants of a cell agree at all four corners and
along all four edges, and differ inside by at most half the cell's
twist,
$\big(\theta_{i,j}+\theta_{i+1,j+1}-\theta_{i,j+1}-\theta_{i+1,j}\big)/2$,
at the center --- zero whenever the four values are locally bilinear.
It is nonetheless a convention, so it is pinned: the triangulation is
computed once, cached, and the same triangles price every iterate.
\end{remark}

\subsection{Beyond the hull: the wing contract}

The grid must end somewhere, and pricing does not.  Below the lowest
strike column, the sheet continues local variance \emph{linearly},
with slope $a$ times the slope of the first interior cell; above the
highest column it is clamped flat.  Three settings of $a$ are
meaningful.  The default, $a=0$, clamps the put wing flat as well.  A
fixed $a>0$ keeps an unquoted put wing rising.  And $a$ is handed to
the optimizer as a free parameter in exactly one circumstance: when a
variance-swap quote is present.  A variance swap pays the realized
variance of the underlying to expiry, and its fair price integrates
the whole smile, deep wings included --- it is the one instrument in
the problem that genuinely reads the deep-put tail, as Chapter~5 will
show.

\begin{remark}[The extrapolation sits outside the positivity guarantee]
\label{rem:lvhull}
\Cref{prop:lvnodal} stops at the hull, and for a reason: outside it,
interpolation becomes extrapolation, and extrapolation weights are
signed --- an extrapolated value is no longer an average of nodal
values.  With $a>0$ and a first cell whose variance decreases toward
the boundary, the linear continuation can reach zero before $y=0$
(pricing floors it there, so the continued curve kinks rather than
turning negative), and a rising continuation can exceed $v_{\rm hi}$
freely.  Every guarantee in this chapter that leans on the box is
therefore scoped to the hull and to the flat-clamped wings.  The wing
is a stated convention --- the same standing as the tail
conventions of Chapters~2--3 --- and the chapter's contract table
records it.
\end{remark}

\begin{remark}[Lee's bound, inherited]
\label{rem:lvlee}
A bounded field bounds the smile, and the bound is worth deriving
because Chapter~3 had to legislate what comes out here for free.  The
first step is that the marched call is \emph{monotone in the field}:
raising local variance anywhere can only raise call prices.  To see
it, let $c_{v}$ and $c_{v+\delta v}$ solve \eqref{eq:lvdupire} for
fields $v$ and $v+\delta v$ with $\delta v\ge0$, and subtract the two
equations.  Their difference $s=c_{v+\delta v}-c_{v}$ starts at
$s(0,\cdot)=0$ and satisfies
\[
\partial_\tau s
=\tfrac12\,(v+\delta v)\,y^{2}\,\partial_{yy}s
+\tfrac12\,\delta v\,y^{2}\,\partial_{yy}c_{v}
\]
(add and subtract $\tfrac12(v+\delta v)y^{2}\partial_{yy}c_{v}$ to
check).  The source term is nonnegative --- $\delta v\ge0$ by
assumption, $\partial_{yy}c_{v}\ge0$ by \cref{prop:lvpropagate} ---
and a heat equation driven by a nonnegative source from a zero start
stays nonnegative (the maximum principle again).  So $s\ge0$.  Now
compare with the constant field $v_{\rm hi}$.  With flat-clamped
wings, $v\le v_{\rm hi}$ holds on the entire strike line, so every
call price is dominated by its price in the Black--Scholes model with
variance rate $v_{\rm hi}$; and since the Black price increases in
total variance (vega is positive), the domination transfers to implied
variance:
\[
w(k,\tau)\;\le\;v_{\rm hi}\,\tau\qquad\text{for every }k .
\]
Total implied variance bounded in strike means its wing slope decays
to zero: the surface sits strictly inside Lee's cap of $2$
automatically, where Chapter~3 had to impose a cap by hand
(\cref{sec:svilee}).  A steep released continuation ($a$ free) escapes
$v_{\rm hi}$, and with it this bound.
\end{remark}
